\section{Star-subalgebras and star-isomorphisms of full matrix algebras}

This section collects identifications of star-subalgebras, the product of two
commuting star-subalgebras, and the unitary implementation of a
star-isomorphism between full matrix algebras on two coordinate sets.

\begin{lemma}[Coordinate identifications preserve ambient values]
    \label{lem:star_subalgebra_coordinate_equivalences}
    \lean{StarSubalgebra.equivMapOfInjective,
        StarSubalgebra.equivMapOfInjective_coe,
        StarSubalgebra.equivOfEq,StarSubalgebra.equivOfEq_coe}
    \leanok
    Let $A,B$ be unital star-algebras over a commutative star-semiring
    $R$. An injective star-algebra homomorphism $f:A\to B$ identifies
    every star-subalgebra $S\subseteq A$ with its image $f(S)$, and the
    identification sends $a\in S$ to the same ambient value $f(a)$.
    Equal star-subalgebras are identified by preserving every ambient
    element. These identifications retain both multiplication and star.
\end{lemma}

\begin{proof}
    \leanok
    Restrict the injective homomorphism to its image, using the usual
    algebra isomorphism, and retain its preservation of star. The equality
    identification changes the membership proof and leaves the element
    itself unchanged, so it also preserves star.
\end{proof}

\begin{definition}[Multiplication of commuting factors]
    \label{def:qca_commuting_factor_multiplication}
    \lean{StarSubalgebra.commutingMulMap,
        StarSubalgebra.commutingMulMap_tmul}
    \leanok
    Let $\mathcal C$ be a nontrivial complex star-algebra and let
    $\mathcal A,\mathcal B\subseteq\mathcal C$ be unital star-subalgebras
    whose elements commute with each other. Multiplication defines a unital
    star-algebra homomorphism
    \begin{align}
        \mu_{\mathcal A,\mathcal B}
        \colon \mathcal A\otimes_{\mathbb C}\mathcal B
         & \longrightarrow \mathcal C, \notag \\
        \mu_{\mathcal A,\mathcal B}(a\otimes b)
         & =ab.
    \end{align}
\end{definition}

\begin{lemma}[Algebraic product of commuting factors]
    \label{thm:qca_commuting_factor_product}
    \lean{StarSubalgebra.commutingMulMap_range,
        StarSubalgebra.sup_toSubmodule_eq_mul}
    \leanok
    \uses{def:qca_commuting_factor_multiplication}
    The image of $\mu_{\mathcal A,\mathcal B}$ is the algebraic join
    $\mathcal A\vee\mathcal B$.  Its underlying complex vector space is
    \begin{align}
        \mathcal A\mathcal B
         & =\operatorname{span}_{\mathbb C}
        \{ab\mid a\in\mathcal A,\ b\in\mathcal B\}.
        \label{eq:qca_commuting_factor_product_span}
    \end{align}
    Thus every element of the join is a finite complex-linear sum of
    products $ab$.
\end{lemma}

\begin{proof}\leanok
    \uses{def:qca_commuting_factor_multiplication}
    The involution reverses products.  Since $a^*\in\mathcal A$ and
    $b^*\in\mathcal B$ commute,
    \begin{align}
        \mu_{\mathcal A,\mathcal B}\bigl((a\otimes b)^*\bigr)
         & =a^*b^*=b^*a^*=(ab)^*.
    \end{align}
    The image contains both factors because $a=a\mathbf 1$ and
    $b=\mathbf 1b$.  Conversely, every pure tensor maps to an element of
    the join, and linearity gives the reverse inclusion.  The description
    of the underlying vector space is the range formula for the linear
    multiplication map.
\end{proof}

\begin{lemma}[Product of two full matrix factors]
    \label{thm:qca_commuting_full_matrix_product}
    \lean{StarSubalgebra.commutingMulMap_injective_of_equiv_matrix,
        StarSubalgebra.commutingSupEquivMatrix,
        StarSubalgebra.nonempty_equiv_sup_matrix,
        StarSubalgebra.finrank_sup_of_equiv_matrix}
    \leanok
    \uses{def:qca_commuting_factor_multiplication,
        thm:qca_commuting_factor_product}
    Suppose that $\mathcal C$ is a nontrivial complex star-algebra,
    that $p,q>0$, and that
    \begin{align}
        \mathcal A & \cong M_p(\mathbb C), \notag \\
        \mathcal B & \cong M_q(\mathbb C)
    \end{align}
    as star-algebras.  Then $\mu_{\mathcal A,\mathcal B}$ is injective and
    \begin{align}
        \mathcal A\vee\mathcal B
         & \cong M_{pq}(\mathbb C), \notag \\
        \dim_{\mathbb C}(\mathcal A\vee\mathcal B)
         & =(pq)^2.
        \label{eq:qca_commuting_full_matrix_product}
    \end{align}
    This is the unital single-block specialization of Schumacher--Werner
    Proposition~\texttt{Cscom}
    \cite[lines~2119--2138]{SchumacherWerner2004QCA}, as used in
    \cite[lines~1270--1308]{Gross2012IndexTheory}; the restriction is
    recorded in \cite{gap:gnvw12_cscom_single_block_scope}.
\end{lemma}

\begin{proof}\leanok
    \uses{def:qca_commuting_factor_multiplication,
        thm:qca_commuting_factor_product}
    The source of the multiplication homomorphism is
    \begin{align}
        M_p(\mathbb C)\otimes_{\mathbb C}M_q(\mathbb C)
         & \cong M_{pq}(\mathbb C).
    \end{align}
    This is a simple ring.  The kernel of a unital homomorphism from a
    simple ring to a nonzero ring is not the whole ring, and is therefore
    zero.  Hence multiplication is injective.  Its image is the join by
    Lemma~\ref{thm:qca_commuting_factor_product}; the matrix presentation
    and dimension formula in
    Equation~(\ref{eq:qca_commuting_full_matrix_product}) follow.
\end{proof}

\begin{lemma}[Prescribed coordinates of commuting products]
    \label{lem:commuting_product_coordinates}
    \lean{StarSubalgebra.commutingMulEquivSup,
        StarSubalgebra.commutingMulEquivSup_tmul_coe,
        StarSubalgebra.commutingSupEquivMatrix_apply_mul}
    \leanok
    \uses{thm:qca_commuting_full_matrix_product}
    For commuting full matrix factors $\mathcal R$ and $\mathcal S$ with
    prescribed presentations $e_R$ and $e_S$, multiplication identifies
    $\mathcal R\otimes\mathcal S$ with $\mathcal R\vee\mathcal S$, and the
    induced matrix coordinates $e_{R\vee S}$ of the join satisfy
    \begin{align}
        e_{R\vee S}(ab) & =e_R(a)\otimes e_S(b).
    \end{align}
\end{lemma}

\begin{proof}
    \leanok
    \uses{thm:qca_commuting_full_matrix_product}
    Multiplication is injective by
    Lemma~\ref{thm:qca_commuting_full_matrix_product} and has the join as its
    range. Transport the tensor product of the two presentations through this
    multiplication isomorphism. A pure tensor maps to its product in the
    ambient algebra and to its Kronecker coordinates.
\end{proof}

\begin{definition}[Conjugation by a rectangular unitary]
    \label{def:rectangular_unitary_conjugation}
    \lean{Matrix.IsUnitaryBetween.conjStarAlgEquiv,
        Matrix.IsUnitaryBetween.conjStarAlgEquiv_apply}
    \leanok
    Let $M,N$ be finite sets and let $U\in\C^{M\times N}$ satisfy
    $U^\dagger U=I_N$ and $UU^\dagger=I_M$. Then $A\mapsto UAU^\dagger$ is a
    star-algebra isomorphism from $M_N(\C)$ to $M_M(\C)$, with inverse
    $B\mapsto U^\dagger BU$.
\end{definition}

\begin{lemma}[A full-matrix star-isomorphism has a unitary implementer]
    \label{lem:qca_full_matrix_spatial_implementer}
    \lean{Matrix.exists_unitary_conj_of_starAlgEquiv_matrix_between}
    \leanok
    \uses{def:rectangular_unitary_conjugation}
    Let $I,J$ be nonempty finite coordinate sets and let
    $f:M_I(\C)\to M_J(\C)$ be a complex star-algebra
    isomorphism. Then $|I|=|J|$ and there is a unitary map
    $U:\C^I\to\C^J$ with
    \begin{align}
        U^\dagger U=\Id_I,\qquad UU^\dagger=\Id_J,\qquad
        f(A)=UAU^\dagger\quad(A\in M_I(\C)).
    \end{align}
    This applies between distinct coordinate sets, as required for the
    unitaries in \cite[equations~\texttt{bruleU} and~\texttt{bruleV},
        lines~980--1001]{SchumacherWerner2004QCA}.
\end{lemma}

\begin{proof}
    \leanok
    The algebra isomorphism is complex linear, so $|I|^2=|J|^2$ and
    $|I|=|J|$. Choose an equivalence of the two coordinate sets and
    transport $f$ to an automorphism of one full matrix algebra.
    The unitary Skolem--Noether theorem implements that automorphism.
    Let $\sigma:J\to I$ be the chosen equivalence, let $P_\sigma$ be its
    permutation matrix, and let $W\in M_I(\C)$ be a unitary with
    $P_\sigma^\dagger f(A)P_\sigma=WAW^\dagger$ for all $A$. Put
    $U=P_\sigma W$. Since multiplication and adjoints respect the
    reindexing,
    \begin{align}
        f(A)
         & =P_\sigma WAW^\dagger P_\sigma^\dagger
        =UAU^\dagger,
         &
        U^\dagger U
         & =W^\dagger W=\Id_I,
         &
        UU^\dagger
         & =P_\sigma P_\sigma^\dagger=\Id_J.
        \label{eq:qca_full_matrix_spatial_transport}
    \end{align}
\end{proof}
